\section{Experiments} \label{sec:05-experiments}
\begin{foldable} % Datasets and models
  We evaluate on two linguistic tasks: AG News ($N=120$\,K), IMDb ($N=25$\,K), SST-2 ($N=67$\,K) for classification; MNLI-m ($N=393$\,K), QNLI ($N=105$\,K), and QQP ($N=364$\,K) for natural language inference (NLI).
  We compare against four baselines:
  \begin{itemize}
    \item \textbf{Random}: uniform random selection at the target budget.
    \item \textbf{EDA}~\cite{wei2019eda}: word-level augmentation (10$\times$) applied to the \textbf{Random} set.
    \item \textbf{DiLM}~\cite{maekawa2024dilm}: LLM-generated candidates selected by gradient matching across k-centre selection. It evaluated on both classification and NLI.
    \item \textbf{DaLLME}~\cite{tao2024textual}: synthetic texts inverted from distilled embeddings. It evaluated on classification only.
  \end{itemize}
  We evaluate TAKE at two budgets (20/cls and 0.1\%) to match DiLM's and DaLLME's reported settings, respectively.
  To match downstream evaluation, downstream models for classification are Logistic Regression~(LR) with TF-IDF~\cite{salton1988term} features, TextCNN~\cite{kim2014convolutional} and TextRNN~\cite{liu2016recurrent} with GloVe~\cite{pennington2014glove} word embeddings, and BERT~\cite{devlin2019bert}.
  For NLI, we evaluate on BERT and Siamese LR.
  Siamese LR is a variant of LR: two input branches to independently encode sentence pairs, and a logistic head over the concatenated representations, tailored as a weak learner for NLI.
  The primary metric is downstream accuracy on classification/NLI tasks, reported as \meanstd{mean}{\text{s.e.}} across five runs.
\end{foldable}

\begin{foldable} % Implementation details
  In terms of implementation, TAKE extracts knowledge scores $\kappa_n$ by running $T=20$ training checkpoints with exponential kernel $\lambda^* = T/T_m$, using all-MiniLM-L6-v2~\cite{wang2020minilm} as the frozen encoder $\phi$.
  The synthetic pool uses Gemma-3~\cite{kamath2025gemma}~(270\,M) fine-tuned on $\mathcal{D}$ to generate $|\mathcal{D}_\text{synth}|=5{,}000$ samples.
  OT selection embeds via the same $\phi$ and uses Sinkhorn with $\varepsilon=0.05$ for 200 iterations.
  All experiments fit on a single NVIDIA V100~(40\,GB).
\end{foldable}

\subsection{Classification} \label{sec:05-results-classification}
\begin{table}[ht]
  \centering
  \begin{threeparttable}
    \caption{Classification accuracy (\%) on AG News, IMDb, and SST-2.}
    \label{tab:results-classification}
    \setlength{\tabcolsep}{4pt}
    \begin{tabular}{l l c c c c c}
      \toprule
      & \textbf{Method} & \textbf{Budget} & \textbf{LR} & \textbf{TextCNN} & \textbf{TextRNN} & \textbf{BERT} \\
      \midrule
      \multirow{5}{*}{\rotatebox{90}{AG News}}
      & Full dataset  & 100\% & \meanstd{90.21}{0.14}    & \meanstd{91.27}{0.20}    & \meanstd{90.98}{0.32}      & \meanstd{93.85}{0.12} \\
      & Random        & 0.1\% & \meanstd{66.39}{2.27}    & \meanstd{75.56}{0.14}    & \meanstd{74.91}{1.45}      & \meanstd{78.12}{3.25} \\
      & EDA           & 0.1\% & \meanstd{68.03}{2.43}    & \meanstd{82.27}{0.15}    & \meanstd{81.46}{1.34}      & \meanstd{84.45}{1.44} \\
      & DaLLME        & 0.1\% & \meanstd{74.60}{0.02}\dg & \meanstd{88.30}{0.02}\dg & \meanstdbf{84.50}{0.02}\dg & $-$ \\
      & \textbf{TAKE} & 0.1\% & \meanstdbf{77.07}{0.09}  & \meanstdbf{88.53}{0.21}  & \meanstd{84.25}{1.03}      & \meanstdbf{89.82}{0.18} \\
      \midrule
      \multirow{5}{*}{\rotatebox{90}{IMDb}}
      & Full dataset  & 100\% & \meanstd{88.41}{0.13}    & \meanstd{87.40}{0.20}    & \meanstd{83.00}{0.35}      & \meanstd{91.22}{0.25} \\
      & Random        & 0.1\% & \meanstd{59.73}{5.51}    & \meanstd{57.65}{1.19}    & \meanstd{55.09}{4.31}      & \meanstd{60.40}{4.10} \\
      & EDA           & 0.1\% & \meanstd{61.52}{4.62}    & \meanstd{62.40}{3.40}    & \meanstd{58.92}{1.93}      & \meanstd{63.12}{1.95} \\
      & DaLLME        & 0.1\% & \meanstd{65.00}{0.02}\dg & \meanstd{68.30}{0.02}\dg & \meanstdbf{64.50}{0.02}\dg & $-$ \\
      & \textbf{TAKE} & 0.1\% & \meanstdbf{66.83}{0.13}  & \meanstdbf{68.93}{2.15}  & \meanstd{63.26}{0.38}      & \meanstdbf{71.55}{0.22} \\
      \midrule
      \multirow{5}{*}{\rotatebox{90}{SST-2}}
      & Full dataset  & 100\%  & \meanstd{80.23}{0.27}   & \meanstd{88.81}{0.15}    & \meanstd{87.40}{0.76}      & \meanstd{92.52}{0.28} \\
      & Random        & 20/cls & \meanstd{55.28}{2.32}   & \meanstd{63.85}{0.87}    & \meanstd{61.52}{1.66}      & \meanstd{69.97}{3.21} \\
      & EDA           & 20/cls & \meanstd{57.21}{1.05}   & \meanstd{65.54}{0.45}    & \meanstd{63.17}{0.84}      & \meanstd{74.21}{1.58} \\
      & DiLM          & 20/cls & $-$                     & $-$                      & $-$                        & \meanstd{80.30}{2.80}\dg \\
      & \textbf{TAKE} & 20/cls & \meanstdbf{65.05}{0.05} & \meanstdbf{70.65}{0.32}  & \meanstdbf{67.18}{0.49}    & \meanstdbf{81.15}{0.24} \\
      \bottomrule
    \end{tabular}
    \begin{tablenotes} \footnotesize
    \item {
        [$-$]~not reported; [$\dagger$]~reported from respective paper; [\textbf{bold}]~best per budget. \phantom{1234}
        DaLLME reported with its best case (OpenAI-3-large embeddings).
      }
    \end{tablenotes}
  \end{threeparttable}
\end{table}

\begin{foldable} % Classification
  \textbf{Accuracy gains and their source.}
  TAKE (0.1\%) outperforms DaLLME on LR by $+2.5$ on AG News and $+1.8$ on IMDb.
  The gap is widest at the LR level and narrows for neural backbones, consistent with TAKE's backbone-agnostic selection---trajectory scores are computed independently of the downstream model.
  TAKE minimizes the OT cost in embedding space, directly controlling the downstream loss gap under $P^w$ via Theorem~\ref{thm:gap} (the gap to the original $P_\mathcal{D}$ is bounded by Remark~\ref{rem:double-gap}).
  Since $\kappa_n$ scores are computed independently of the evaluation backbone, this coverage of $P^w$ transfers across model families.
\end{foldable}

\begin{foldable}
  \textbf{Stability across runs.}
  The wide standard errors for Random (IMDb LR $\pm5.5$, BERT $\pm4.1$) reveal that random selection at 0.1\% can land anywhere from a coherent corpus to a degenerate one.
  TAKE's errors are at least $5\times$ narrower ($\pm0.1$--$2.2$), a direct consequence of OT-based selection: the transport plan distributes selected prototypes across $P^w$ rather than drawing them independently, suppressing the variance that Random incurs by chance.
\end{foldable}

\begin{foldable}
  \textbf{SST-2 vs.\ DiLM.}
  At 20/cls, TAKE matches DiLM on BERT ($81.15$ vs.\ $80.30$, within s.e.) while also providing results for LR, TextCNN, and TextRNN where DiLM has none---indicating that the OT selection step generalizes across backbone families without retraining the scorer.
\end{foldable}

\subsection{Natural Language Inference} \label{sec:05-results-nli}
\begin{table}[ht]
  \centering
  \begin{threeparttable}
    \caption{NLI accuracy (\%) on MNLI-m and QQP.}
    \label{tab:results-nli}
    \setlength{\tabcolsep}{4pt}
    \begin{tabular}{l l c c c}
      \toprule
      & \textbf{Method} & \textbf{Budget} & \textbf{Siamese LR} & \textbf{BERT} \\
      \midrule
      \multirow{5}{*}{\rotatebox{90}{MNLI-m}}
      & Full dataset    & 100\%  & \meanstd{60.02}{0.07}   & \meanstd{86.81}{0.30} \\
      & Random          & 20/cls & \meanstd{34.15}{2.10}   & \meanstd{40.10}{3.20} \\
      & EDA             & 20/cls & \meanstd{37.80}{1.45}   & \meanstd{44.52}{2.80} \\
      & DiLM            & 20/cls & $-$                     & \meanstd{48.70}{2.60}\dg \\
      & \textbf{TAKE}   & 20/cls & \meanstdbf{42.66}{0.15} & \meanstdbf{51.24}{0.45} \\
      \midrule
      \multirow{5}{*}{\rotatebox{90}{QNLI}}
      & Full dataset    & 100\%  & \meanstd{73.05}{0.13}   & \meanstd{82.50}{0.12} \\
      & Random          & 20/cls & \meanstd{52.21}{0.28}   & \meanstd{57.19}{2.20} \\
      & EDA             & 20/cls & \meanstd{54.33}{1.21}   & \meanstd{59.81}{1.42} \\
      & DiLM            & 20/cls & $-$                     & $-$ \\
      & \textbf{TAKE}   & 20/cls & \meanstdbf{57.96}{0.58} & \meanstdbf{64.89}{3.58} \\
      \midrule
      \multirow{5}{*}{\rotatebox{90}{QQP}}
      & Full dataset    & 100\%  & \meanstd{78.58}{0.05}   & \meanstd{89.45}{0.15} \\
      & Random          & 20/cls & \meanstd{52.40}{1.32}   & \meanstd{59.10}{3.80} \\
      & EDA             & 20/cls & \meanstd{55.12}{1.88}   & \meanstd{62.35}{2.45} \\
      & DiLM            & 20/cls & $-$                     & \meanstd{64.40}{2.20}\dg \\
      & \textbf{TAKE}   & 20/cls & \meanstdbf{60.04}{0.12} & \meanstdbf{67.76}{0.35} \\
      \bottomrule
    \end{tabular}
    \begin{tablenotes} \footnotesize
    \item {
        [$-$]~not reported; [$\dagger$]~reported from respective paper;
        [\textbf{bold}]~best per budget.
      }
    \end{tablenotes}
  \end{threeparttable}
\end{table}
\begin{foldable} % NLI
  \textbf{Gains and their interpretation.}
  DaLLME has no published NLI experiments and is excluded from this comparison.
  In comparison with DiLM, TAKE outperforms DiLM by $+2.5$ on MNLI-m and $+3.4$ on QQP on a BERT learner, and provides competitive results for QNLI.
  The LR gap ($+8.5$ on MNLI-m, $+5.8$ on QNLI, $+7.6$ on QQP relative to Random) is particularly informative: Siamese LR has no capacity to compensate for a poor corpus---every percentage point of accuracy here must come from the quality of $\tilde{\mathcal{D}}$ itself.
  That TAKE's Siamese LR exceeds the Random BERT baseline on both QQP ($60.04$ vs.\ $59.10$) and QNLI ($57.96$ vs.\ $57.19$) suggests the selected sentence pairs are genuinely more informative, not merely better distributed.
\end{foldable}

\begin{foldable}
  \textbf{The NLI ceiling gap and what it implies.}
  All distilled methods fall short of the full-dataset ceiling ($86.8$ / $82.5$ / $89.5$ for MNLI-m, QNLI, QQP), a gap larger than on single-sentence tasks.
  NLI imposes a joint constraint: the synthetic pool must contain pairs whose relationship correctly reflects each class label at extreme budget ($\approx\!60$ pairs for MNLI-m, $\approx\!40$ for QQP and QNLI, at 20/cls).
  This is a harder coverage requirement than single-sentence classification, and the synthetic pool must cover all relationship types within it.
  This suggests that pool generation and budget are the binding constraints, rather than scoring or selection.
\end{foldable}

\begin{foldable}
  \textbf{Stability.}
  On MNLI-m and QQP, TAKE's errors ($\pm0.1$--$0.5$) are $5$--$7\times$ narrower than DiLM's ($\pm2.2$--$2.6$), mirroring the pattern on classification.
  This is consistent with the OT transport plan enforcing coverage of $P^w$ rather than selecting samples independently: deterministic coverage implies that any two runs of TAKE over the same pool yield nearly identical corpora, whereas DiLM's gradient-matching step introduces variance through its dependence on a randomly initialized model.
\end{foldable}

\subsection{Component Contributions} \label{sec:05-results-components}
\begin{table}[ht]
  \centering
  \begin{threeparttable}
    \caption{Distillation accuracy (\%) under TAKE component ablation.}
    \label{tab:results-abla-components}
    \setlength{\tabcolsep}{4pt}
    \begin{tabular}{l c c c c}
      \toprule
      & \multicolumn{2}{c}{\textbf{AG News}} & \multicolumn{2}{c}{\textbf{QQP}} \\
      \cmidrule(lr){2-3} \cmidrule(lr){4-5}
      & \textbf{LR} & \textbf{BERT} & \textbf{Siamese LR} & \textbf{BERT} \\
      \midrule
      Random                   & \meanstd{66.39}{2.27}   & \meanstd{75.92}{0.98}   & \meanstd{52.40}{1.32}   & \meanstd{57.94}{1.30} \\
      \bltrg $k$-Means nearest & \meanstd{72.89}{1.37}   & \meanstd{77.79}{1.04}   & \meanstd{55.28}{1.12}   & \meanstd{61.30}{1.08} \\
      TAKE without $\kappa$    & \meanstd{75.82}{0.23}   & \meanstd{80.94}{0.93}   & \meanstd{58.80}{0.92}   & \meanstd{63.66}{1.25} \\
      \textbf{TAKE}            & \meanstdbf{77.07}{0.09} & \meanstdbf{86.50}{1.01} & \meanstdbf{60.04}{0.12} & \meanstdbf{66.29}{1.91} \\
      \bottomrule
    \end{tabular}
    \begin{tablenotes} \footnotesize
      \item {[$\blacktriangle$]~\textbf{not} inclusive of TAKE; [\textbf{bold}]~best results.}
    \end{tablenotes}
  \end{threeparttable}
\end{table}

\begin{foldable} % Ablation
  To isolate the contribution of each design choice, we ablate TAKE's components on AG News and QQP as representatives for the classification and NLI tasks, in Table~\ref{tab:results-abla-components}.
  TAKE improves over Random by $+10.7$ on AG News LR and $+10.6$ on BERT; $k$-means, by comparison, accounts for only part of this gain ($+6.5$ LR, $+1.9$ BERT), while our discrete OT selection step alone (TAKE without $\kappa$) already achieves $+9.4$ LR and $+5.0$ BERT from the same baseline.
  The same pattern holds on QQP: $k$-means gains $+2.9$ / $+3.4$ (Siamese LR / BERT) whereas TAKE without $\kappa$ gains $+6.4$ / $+5.7$.
  Adding $\kappa$ knowledge scores delivers the remaining improvement, modestly for Siamese LR ($+1.2$) but substantially for BERT ($+5.6$ on AG News, $+2.6$ on QQP), suggesting that $\kappa$ primarily benefits models with enough capacity to exploit the harder, informative examples that high-$\kappa$ samples represent.
\end{foldable}

\subsection{Kernels for Temporal Weighting} \label{sec:05-results-kernels}
\begin{table}[ht]
  \centering
  \begin{threeparttable}
    \caption{Distillation accuracy (\%) with alternative temporal kernels.}
    \label{tab:results-abla-kernels}
    \setlength{\tabcolsep}{4pt}
    \begin{tabular}{l c c c c}
      \toprule
      & \multicolumn{2}{c}{\textbf{AG News}} & \multicolumn{2}{c}{\textbf{QQP}} \\
      \cmidrule(lr){2-3} \cmidrule(lr){4-5}
      \textbf{Kernel} & \textbf{LR} & \textbf{BERT} & \textbf{Siamese LR} & \textbf{BERT} \\
      \midrule
      TAKE without $\kappa$       & \meanstd{75.82}{0.23}   & \meanstd{80.94}{0.93}   & \meanstd{58.80}{0.92}   & \meanstd{63.66}{1.25} \\
      Last checkpoint             & \meanstd{75.24}{0.57}   & \meanstd{78.28}{0.66}   & \meanstd{58.32}{1.56}   & \meanstd{62.19}{0.57} \\
      Constant                    & \meanstd{76.13}{1.14}   & \meanstd{83.47}{0.71}   & \meanstd{59.44}{0.86}   & \meanstd{64.78}{1.51} \\
      Linear                      & \meanstd{76.20}{0.68}   & \meanstd{85.20}{1.00}   & \meanstd{59.61}{1.11}   & \meanstd{65.60}{1.35} \\
      Cosine                      & \meanstd{76.74}{0.53}   & \meanstd{85.46}{0.64}   & \meanstd{59.14}{0.85}   & \meanstd{66.16}{0.41} \\
      Exponential (\textbf{TAKE}) & \meanstdbf{77.07}{0.09} & \meanstdbf{86.50}{1.01} & \meanstdbf{60.04}{0.12} & \meanstdbf{66.29}{1.91} \\
      \bottomrule
    \end{tabular}
    \begin{tablenotes} \footnotesize
    \item {[\textbf{bold}]~best results.}
    \end{tablenotes}
  \end{threeparttable}
\end{table}

\begin{foldable}
  Table~\ref{tab:results-abla-kernels} ablates the choice of temporal kernel defined in \S\ref{sec:04-gradient-kappa}.
  The last-checkpoint baseline is the critical exception: it falls \textit{below} the no-$\kappa$ baseline across all backbones, most severely on BERT ($-2.7$ on AG News, $-1.5$ on QQP).
  This is precisely the failure mode identified in Proposition~\ref{prop:traj}: when $\kappa$ is extracted at convergence only, the learner has fully memorized difficult samples, so the \textit{hard-sample bias} is intensified rather than corrected, and the selected corpus is skewed towards harder, less representative examples.
  All other trajectory-based strategies lessen this effect and outperform the no-$\kappa$ baseline (OT only), with decreasing kernels (linear, cosine, exponential) doing so most strongly by concentrating mass on the early-to-mid training phase where the bias-correcting signal is strongest.
  The exponential kernel (TAKE default) achieves the best result in all cases, consistent with the theoretical preference for early-phase concentration of influence.
\end{foldable}
